\section{Bell's theorem and the retrocausal worldtube interpretation}
\label{sec:bell-interpretation}

The theory developed in \Cref{sec:continuum-model,sec:emergent-relativity,sec:geophase,sec:effective-field,sec:localized} is manifestly local and deterministic: a classical brane action whose dynamics is fixed by the embedded configuration, with all emergent objects (effective metric, gauge connections, soliton labels) derived as read-outs of the substrate state.
Bell's theorem then forces an interpretive stance.
The 4D-in-4D ontology of \Cref{subsec:ontology-kinematics} commits the present program to a specific resolution, which we make explicit in this section.

\subsection{The Bell constraint}

Bell's theorem rules out theories that are simultaneously
(i) \emph{local} in the sense that no influence propagates outside the future lightcone of an event,
(ii) \emph{deterministic} in the sense that hidden variables fully fix outcomes given the measurement context,
and (iii) \emph{measurement-independent} in the sense that the joint distribution of detector settings and hidden state factorises over distant detectors \citep{Bell1964,CHSH1969}.
The experimental violation of Bell inequalities, by now closing the standard loopholes \citep{Hensen2015,Giustina2015,Shalm2015}, then implies that at least one of (i)--(iii) must be given up by any classical-substrate completion that aims to reproduce the empirical correlations of entangled systems.

A literal superdeterministic completion --- giving up (iii) by correlating detector settings with the hidden state at the Big Bang --- is logically consistent but is generally regarded as explanatorily empty: it postpones rather than resolves the question of where correlations come from \citep{Hossenfelder2020Superdeterminism}.

\subsection{The retrocausal loophole and where this theory lives}

The alternative consistent with the substrate ontology of \Cref{subsec:ontology-kinematics} is the \emph{retrocausal} loophole, which relaxes (iii) not by tying detector settings to a cosmic initial condition but by allowing the future measurement context to propagate back along the particle's worldtube to its origin in the shared past of the entangled pair.
Spacelike correlations are then not action-at-a-distance but \emph{continuity of a single extended 4D object} whose world-volume branches at a V-vertex in the past.
The present theory places itself in the same interpretive family as the two-state-vector formalism of Aharonov and collaborators \citep{AharonovBergmannLebowitz1964,Aharonov2008TSVF}, Cramer's transactional interpretation \citep{Cramer1986TI}, the retrocausal models of Price and Wharton \citep{PriceWharton2015Retrocausality}, and Sutherland's time-symmetric Bohmian model \citep{Sutherland2017}.

This stance fits the 4D-in-4D ontology naturally.
The brane action is time-symmetric: its Lorentzian sign structure (\Cref{subsec:ontology-kinematics}) picks out a \emph{kind} of direction (timelike vs spacelike) but not a \emph{direction} along the timelike axis (past-to-future vs future-to-past).
The arrow of time is therefore not a property of the substrate; it is selected at the soliton level by the chirality of soliton solutions.

\subsection{Causality from soliton chirality; matter and antimatter}
\label{subsec:soliton-chirality-bell}

In this stance, the empirical arrow of time is a property of \emph{solitons}, not of the substrate.
Matter solitons are forward-propagating worldtubes; antimatter solitons are backward-propagating worldtubes.
This realises the Feynman--Stueckelberg interpretation of antiparticles as time-reversed particles \citep{Stueckelberg1941,Feynman1949ApproachQED} at the level of substrate soliton solutions rather than as an interpretive remark on a Lorentz-invariant field theory.
The cosmological baryon--antibaryon asymmetry is, on this reading, a \emph{geometric initial-condition asymmetry}: matter and antimatter solitons emerge from the Big Bang vertex propagating in opposite time directions along the brane's selected timelike axis, and our universe contains only the forward-propagating branch.
This differs structurally from Sakharov-condition baryogenesis \citep{Sakharov1967CPviolation} in which the same time direction hosts both matter and antimatter, with a small dynamical asymmetry favouring matter.

\subsection{Entanglement as V-branching worldtubes}

An entangled pair of solitons, on the present reading, is \emph{one} branching 4D world-volume whose worldtube splits at a V-vertex in the shared past.
Each branch propagates forward to a measurement event; the boundary condition imposed by the measurement on one branch is communicated to the other branch via continuity through the shared V-vertex.
The result is statistical correlation that is local along each worldline (no superluminal influence on any spacelike slice the lab observer can resolve) and deterministic given the full 4D configuration, while still violating Bell inequalities because the joint measurement context affects the past of both branches via the V-vertex.
The 4D world-volume picture promotes ``spooky action at a distance'' to ``ordinary continuity of one extended object''.

\subsection{Open derivations}
\label{subsec:bell-open}

The interpretive stance above places three concrete derivational obligations on the program.
None of these is claimed as established in this paper; they are listed as honest debts of the framework and as the constraints on what the theory can responsibly assert about quantum foundations.

\paragraph{(i) Tsirelson's bound.}
Retrocausal models notoriously have difficulty reproducing the exact Tsirelson bound $2\sqrt{2}$ for the maximum CHSH correlation \citep{Cirelson1980}: they tend to under- or over-shoot.
A successful completion of the present theory requires deriving this bound from the soliton-correlation structure of the V-branching worldtube, not merely allowing nonzero correlation.
This is the most direct empirical handle on the stance.

\paragraph{(ii) No-signalling theorem.}
In standard quantum mechanics, no-signalling is automatic from the partial-trace structure of bipartite states.
In the present substrate it must be derived as a theorem about the dynamics of branching worldtubes --- specifically, that no statistical observable accessible on one branch depends on the measurement context of the other branch.
This is the protection that prevents the retrocausal channel from being used to send a bit backward in time at the macroscopic level.

\paragraph{(iii) Baryon-to-photon ratio.}
The geometric-asymmetry picture of matter and antimatter expanding from the Big Bang vertex in opposite time directions is appealing as a unification of baryogenesis with the Bell stance.
However, the observed baryon-to-photon ratio $\eta\simeq 6\times 10^{-10}$ is currently fit by Sakharov-condition models with several parameters \citep{Sakharov1967CPviolation,Riotto1999BaryogenesisReview}.
Whether the geometric-vertex picture can match this number quantitatively, rather than just allow opposite-direction expansion in principle, is an open question.
